\section*{第 \ref{ch:g6:density-floating} 章——密度：物体为何漂浮}

\begin{solution}{exo:g6:density-floating:1}
该物质每\omterm{def:g6:mass-and-volume:volume}{立方厘米}的\omterm{def:g3:measuring-mass:mass}{质量}，以克计；
水的数值为 \qty{1}{g} 每 \unit{cm^3}。规律：
比水\omterm{def:g6:density-floating:density}{密度}小的\omterm{def:g1:floating-and-sinking:words}{漂浮}，比水大的\omterm{def:g1:floating-and-sinking:words}{下沉}。
\end{solution}

\begin{solution}{exo:g6:density-floating:2}
橡木\omterm{def:g1:floating-and-sinking:words}{漂浮}；石头\omterm{def:g1:floating-and-sinking:words}{下沉}；软木\omterm{def:g1:floating-and-sinking:words}{漂浮}；铁\omterm{def:g1:floating-and-sinking:words}{下沉}；冰\omterm{def:g1:floating-and-sinking:words}{漂浮}——
仅差毫厘。
\end{solution}

\begin{solution}{exo:g6:density-floating:3}
$120 \div 80 = 1.5$：\omterm{def:g6:density-floating:density}{密度}为 $1.5$——比水大，所以
\omterm{def:g1:floating-and-sinking:words}{下沉}。
\end{solution}

\begin{solution}{exo:g6:density-floating:4}
$158 \div 20 = 7.9$：这是铁的特征值。
\end{solution}

\begin{solution}{exo:g6:density-floating:5}
\omterm{def:g6:density-floating:density}{密度}是物质的属性，而非某一块的属性：每\omterm{def:g6:mass-and-volume:volume}{立方厘米}
橡木所含克数相同，无论碎屑还是
大梁。因此\omterm{def:g6:density-floating:density}{密度}可作为未知物质的“身份证”。
\end{solution}

\begin{solution}{exo:g6:density-floating:6}
\omterm{def:g3:solids-liquids-gases:liquid}{液体}按\omterm{def:g6:density-floating:density}{密度}分层，\omterm{def:g6:density-floating:density}{密度}最大者居下：蜂蜜 $1.4$
在下，水 $1.0$ 居中，油 $0.9$ 在最上。葡萄（$1.05$）\omterm{def:g1:floating-and-sinking:words}{下沉}
穿过油与水，停在蜂蜜这一层——
比上面两层都重，却比蜂蜜轻。
\end{solution}

\begin{solution}{exo:g6:density-floating:7}
多数\omterm{def:g3:solids-liquids-gases:solid}{固体}比同种\omterm{def:g3:solids-liquids-gases:liquid}{液体}\omterm{def:g6:density-floating:density}{密度}大，在其中\omterm{def:g1:floating-and-sinking:words}{下沉}。冰
能\omterm{def:g1:floating-and-sinking:words}{漂浮}，是因为水\omterm{def:g6:water-states:changes}{凝固}时\omterm{def:g6:mass-and-volume:volume}{体积}膨胀——同样的物质
占据多十分之一的空间，\omterm{def:g6:density-floating:density}{密度}为 $0.9 < 1$。
\end{solution}

\begin{solution}{exo:g6:density-floating:8}
木材\omterm{def:g6:density-floating:density}{密度}约 $0.6$：原木每\omterm{def:g6:mass-and-volume:volume}{立方厘米}
都轻于水的 \qty{1}{g}，故无论多大都会\omterm{def:g1:floating-and-sinking:words}{漂浮}。
青铜\omterm{def:g6:density-floating:density}{密度}约 $9$：硬币每\omterm{def:g6:mass-and-volume:volume}{立方厘米}
都重于水，故无论多小都会\omterm{def:g1:floating-and-sinking:words}{下沉}。大小从来不是关键；
真正起决定作用的始终是\omterm{def:g6:measurement-in-science:measuring}{单位}\omterm{def:g6:mass-and-volume:volume}{体积}的\omterm{def:g3:measuring-mass:mass}{质量}。
\end{solution}

\begin{solution}{exo:g6:density-floating:9}
$\qty{1}{L} = \qty{1000}{cm^3}$，故 $800 \div 1000 = 0.8$
克每\omterm{def:g6:mass-and-volume:volume}{立方厘米}。冰的\omterm{def:g6:density-floating:density}{密度} $0.9$ 比此\omterm{def:g3:solids-liquids-gases:liquid}{液体}
\emph{大}，冰块在其中\omterm{def:g1:floating-and-sinking:words}{下沉}。
\end{solution}

\begin{solution}{exo:g6:density-floating:10}
\omterm{def:g1:floating-and-sinking:words}{漂浮}的整体是整个船体：钢壳\emph{加上}
其所围住的大量\omterm{def:g2:air-around-us:air}{空气}。若按全部空间平均，
整体\omterm{def:g6:density-floating:density}{密度}降到 $1$ 以下，规律便说会\omterm{def:g1:floating-and-sinking:words}{漂浮}。若
船体出现破洞，水取代\omterm{def:g2:air-around-us:air}{空气}，整体\omterm{def:g6:density-floating:density}{密度}上升
趋近钢材本身的\omterm{def:g6:density-floating:density}{密度}，超过 $1$——而规律不变，
则会\omterm{def:g1:floating-and-sinking:words}{下沉}。
\end{solution}

\begin{solution}{exo:g6:density-floating:11}
标准换了：相对于海水的 $1.03$，游泳者
\omterm{def:g6:density-floating:density}{密度}接近 $1$ 的身体明显位于\omterm{def:g1:floating-and-sinking:words}{漂浮}一侧，所以海水
托你浮得更高。进入淡水（$1.0$）时，船失去
这一余量，\omterm{def:g1:floating-and-sinking:words}{下沉}更深：水线沿船身上升。
\end{solution}

\begin{solution}{exo:g6:density-floating:12}
半浸时，球体在水中的深度恰与
一个\omterm{def:g6:density-floating:density}{密度}应为水一半的整体相符：球壳加\omterm{def:g2:air-around-us:air}{空气}的
平均\omterm{def:g6:density-floating:density}{密度}约为 $0.5$ 克每\omterm{def:g6:mass-and-volume:volume}{立方厘米}——隐藏的\omterm{def:g2:air-around-us:air}{空气}
承担了大部分的减重作用。水渗入时取代\omterm{def:g2:air-around-us:air}{空气}，整体
平均\omterm{def:g6:density-floating:density}{密度}上升，球体\omterm{def:g1:floating-and-sinking:words}{下沉}越来越深，
直到平均\omterm{def:g6:density-floating:density}{密度}超过水的 $1.0$ 那一刻，
\omterm{def:g1:floating-and-sinking:words}{漂浮}规律翻转，球体便沉下去。
\end{solution}

\begin{solution}{pb:g6:density-floating:1}
\textbf{1.} 没有变化：用等\emph{\omterm{def:g3:measuring-mass:mass}{质量}}的银替换金，
\omterm{def:g3:levers-and-scales:scale}{天平}仍然平衡——单凭\omterm{def:g3:measuring-mass:mass}{质量}看不出
这种替换。
\textbf{2.} $1930 \div 19.3 = \qty{100}{cm^3}$。
\textbf{3.} \qty{130}{cm^3}——溢出的水量即皇冠的
\omterm{def:g6:mass-and-volume:volume}{体积}，依据排水法得出。
\textbf{4.} 皇冠比同\omterm{def:g3:measuring-mass:mass}{质量}的纯金多占 \qty{30}{cm^3}
的\omterm{def:g6:mass-and-volume:volume}{体积}，说明内部藏有某种
比金更蓬松的物质。
\textbf{5.} $1930 \div 130 \approx 14.8$。
\textbf{6.} 金 $19.3$——皇冠 $14.8$——银 $10.5$：皇冠的
\omterm{def:g6:density-floating:density}{密度}落在两种特征值之间。
\textbf{7.} “陛下，每\omterm{def:g6:mass-and-volume:volume}{立方厘米}纯金的\omterm{def:g3:measuring-mass:mass}{质量}为
\qty{19.3}{g}；您的皇冠仅约 $15$，
并非纯金。”
\textbf{8.} 误差范围给出的\omterm{def:g6:density-floating:density}{密度}为 $1900 \div 130 \approx
14.6$ 至 $1960 \div 130 \approx 15.1$——远非
$19.3$：任何合理的\omterm{def:g3:levers-and-scales:scale}{天平}误差都无法为金匠开脱。
\textbf{9.} 银的\omterm{def:g6:density-floating:density}{密度}较小：每偷走一克金，
便换回一克需\emph{占据更大空间}的银——
等克数替换保持了\omterm{def:g3:measuring-mass:mass}{质量}，却使\omterm{def:g6:mass-and-volume:volume}{体积}膨胀。
\textbf{10.} 金的部分：$965 \div 19.3 = \qty{50.0}{cm^3}$；
银的部分：$965 \div 10.5 \approx \qty{91.9}{cm^3}$。
\textbf{11.} 合计 $\approx \qty{142}{cm^3}$——多于
实测的 $130$：对半猜测用掉的银过多；
真实的窃取量应更少（实测\omterm{def:g6:mass-and-volume:volume}{体积}位于
$100$ 与 $142$ 之间）。
\textbf{12.} “已证明：皇冠并非纯金。方法：其
\omterm{def:g3:measuring-mass:mass}{质量}除以用水测得的\omterm{def:g6:mass-and-volume:volume}{体积}，得出的\omterm{def:g6:density-floating:density}{密度}
远低于金不变的特征值。况且皇冠
连一道划痕都不需要：水从外部
读出了它的秘密——这正是胜利所在。”
\end{solution}
